\section{Recall and settings}
\subsection{Measuring recall}
The methods are evaluated against the top results obtained by scoring every document. These are the \emph{reference results}.
Let $G_K(q)$ be that list of top-$K$ IDs and $A_K(q)$ the IDs returned by the method being tested.
For $K\le N$, recall for query $q$ is the fraction of reference IDs we recover:
\begin{equation}
R_K(q)=\frac{|G_K(q)\cap A_K(q)|}{K}.
\label{eq:recall}
\end{equation}
The intersection symbol means IDs present in both lists. In our example, $G_2=[3,4]$.
Returning [3,0] recovers one of the two, giving 50\% recall. Returning [3,4] gives 100\%.
Finding many candidates is not enough: the correct IDs must be among them and remain in the
final result list.

Here, ``correct'' means the ranking from the binary-document score in
Equation~\ref{eq:score}. Comparing instead with full float vectors would also count changes
caused by quantization. Human relevance labels ask whether a document's text is useful for
the query. These are different tests: matching the binary ranking does not by itself prove
that the text is relevant or that a language model will answer correctly.

If we follow a reference document through the pipeline, we can lose it when its cluster is not
opened, or when local search misses it inside an opened cluster. Fix a query and choose one
of its $K$ reference IDs uniformly. Let $V$ mean that we visit its cluster, and $L$ mean that
the local method returns it. Then
\begin{equation}
R_K(q)=\Pr(V\mid q)\Pr(L\mid V,q).
\label{eq:survival}
\end{equation}
The first probability asks how often a reference document's cluster is opened. The second
asks how often it is returned \emph{given that its cluster was opened}. For example, suppose
8 of the correct 10 documents reach the local search and it returns 6 of those 8. Overall
recall is $(8/10)(6/8)=6/10$. We do not need to assume those two stages are independent.

The scan returns every reference result that reaches it, so its second probability is one.
B and C can also miss documents inside an opened cluster if we stop them early. Across many
queries, average the complete per-query recall; multiplying two separately averaged stage
probabilities generally gives a different answer.

\subsection{A recall model}
We can calculate an average recall if we first specify how the example documents are generated.
Suppose each document bit is chosen independently, with equal chance of 0 or 1. Let $m$ query coordinates have a large magnitude $a$, and let the remaining $d-m$ coordinates
have a small magnitude $\epsilon$, with
\begin{equation}
a>(d-m)\epsilon.
\label{eq:dominance}
\end{equation}
The left side is the weight of one important coordinate; the right side is the total weight
of all the small coordinates. This condition makes one important mismatch more costly than
all the small mismatches together. Therefore a document matching every important coordinate
beats any document that misses one. This is an assumption we choose for the example, not a
property promised by Matryoshka training.

Each document has probability $p=2^{-m}$ of matching all $m$ important signs. Across $N$
independently generated documents, let $Z$ be the number that match. This count follows a
\emph{binomial distribution}: a formula for how many successes occur when we repeat the same
independent chance of success. It gives
\[
\Pr(Z=z)=\binom Nz p^z(1-p)^{N-z},\qquad \E[Z]=Np.
\]
Here $\Pr(Z=z)$ is the chance of exactly $z$ matching documents. The symbol $\binom Nz$
counts the ways those $z$ documents can be chosen. The expression $\E[Z]$ means the average
count across repeated random collections; it is not a guaranteed count for one collection.

If we open only the matching group and score all its rows, we find the whole top-$K$ when
$Z\ge K$. When $Z<K$, we find only $Z$ of them. Recall is therefore $\min(Z,K)/K$.
Multiply each possible recall by its probability and add:
\begin{equation}
\E[R_K]=\frac1K\sum_{z=0}^{N}\min(z,K)\binom Nz p^z(1-p)^{N-z}
       =\frac1K\sum_{i=1}^{K}\Pr(Z\ge i).
\label{eq:binomial-recall}
\end{equation}
The second form is the same average written another way: add the chances of finding at least
one match, at least two, and so on through $K$. It avoids adding a term for every possible
number of documents. The example below shows the calculation with just two documents.

Consider $N=2$, $m=1$, and $K=1$. The ideal group is empty with probability $1/4$, contains
one document with probability $1/2$, and contains two with probability $1/4$. We can calculate
expected recall directly: $0(1/4)+1(1/2)+1(1/4)=3/4$. Notice that the average group size is
one, but some groups are empty. If we substituted that average as though it were a guaranteed
count, we would incorrectly predict 100\% recall.

\fig{recall-example}{Two independent documents and one required matching bit. Opening only the matching group succeeds with probability $3/4$. The empty-group outcome explains why the mean group size is insufficient.}{.84}

At $N=10^6$ and $m=13$, the ideal group averages about 122 rows. With $K=100$, some groups
can still be too small, so a complete B or C search may need to continue into other groups.
Equation~\ref{eq:binomial-recall} describes stopping after the ideal group; it does not describe
every search policy. If we increase a work budget, we have to consider which additional groups
we visit and the scores and sizes of those groups.

For real queries, keeping half the coordinates does not mean keeping half the information
needed to rank documents. Some bits may tend to change together, query weights differ, and
cluster selection may miss good documents. All of these affect recall. An earlier estimate
used a normal distribution---a probability model centered on an average---to approximate
prefix and full scores. It was fitted to some
queries and checked on others, but the Nomic check missed recall by as much as about 9.6
percentage points. That failed the proposed two-point limit, so we cannot use that estimate
to claim real-data recall at a billion rows.

\subsection{Choosing settings}
Some choices change what we build: the code width $d$, cluster count $C$, routing method,
and C's key coordinates and width $h$. Others can change from query to query: probes $P$,
B's leaf size and node budget, C's lookup and candidate limits, and the requested $K$. If
we change $d$, the reference ranking changes too, so we need to compare all methods again
using that same representation.

We can use the cost equation to narrow the settings worth trying. Suppose B follows one
promising path, each split halves the candidates, and the masks keep width $W=\ceil{N/w}$.
After $j$ splits, we expect $N2^{-j}$ rows to reach the scorer. If we temporarily leave out
changes in queue cost, this gives
\begin{equation}
T_B(j)\approx T_0+jWc_{\rm split}+Wc_{\rm leaf}+N2^{-j}c_{\rm score}.
\label{eq:depth}
\end{equation}
If we make one more split, we save about $N2^{-(j+1)}$ scores and pay for another full-width
split. It helps in this model when the saved scoring time exceeds $Wc_{\rm split}$ and the
additional queue cost. For example, with $N=6{,}400$ and $j=5$, another balanced cut saves
about 100 scores. We compare the time of those 100 scores with the time of a 100-word split;
the two counts alone are not comparable.

The number of splits must be a whole number, but we can first treat $j$ as a continuous
variable to find where this time formula stops decreasing. The derivative measures that
change in estimated time. Setting it to zero gives a candidate depth:
\[
\frac{dT_B}{dj}=Wc_{\rm split}-N\ln(2)2^{-j}c_{\rm score}=0,
\qquad
j^*=\log_2\!\left(\frac{N\ln(2)c_{\rm score}}{Wc_{\rm split}}\right).
\]
If the formula suggested 3.6 splits, for example, we would test the valid whole-number choices
3 and 4. We then check the score assumptions and required recall. The
calculation assumes one useful balanced path;
it does not cover arbitrary branch exploration. Since $W$ grows with $N$, a larger
collection does not automatically call for a larger optimal depth. The dense masks themselves
still take more work as the collection grows.

The document counts also help us choose when to stop splitting. At one million rows, the
13-bit matching group contains about 122 rows on average. Its standard deviation is about 11;
this measures how much the group count varies around the average, not a limit that every
group must obey. A threshold of 128 can make a slightly larger group split again on an
unimportant coordinate. The rough margin $122+3(11)\approx155$ suggested trying $L=160$.
The preceding 12-bit group averages about 244 rows, so it will usually still split. We tested
this setting; the rough margin alone does not prove it is best.

The complete choice for method $a$ is
\begin{equation}
\min_{\theta\in\Theta_a}T_a(\theta)
\quad\text{subject to}\quad R_a(\theta)\ge r,\quad M_a(\theta)\le M_{\rm limit}.
\label{eq:choice}
\end{equation}
In this formula, $\theta$ is a set of settings, such as cluster count and leaf size.
We want the smallest query time among settings that meet both the recall target and the RAM
limit. A derivative can help with a known time formula, but $N$ and RAM do not tell us how
many documents a real cluster contains or how often routing misses a good result. We still
need measurements. We compare the best \emph{tested} choices, including good routing options
for A and a large enough leaf for B to scan when further splitting is not useful.